\documentclass[11pt]{article}

\usepackage[a4paper,margin=27mm]{geometry}
\usepackage[T1]{fontenc}
\usepackage{lmodern}
\usepackage{microtype}
\usepackage{amsmath,amssymb,amsthm,mathtools}
\usepackage{array,booktabs,tabularx}
\usepackage{enumitem}
\usepackage{xurl}
\usepackage[hidelinks]{hyperref}
\setlength{\emergencystretch}{3em}
\newcolumntype{Y}{>{\raggedright\arraybackslash}X}
\newtheorem{proposition}{Proposition}

\title{\textbf{Experimental protocol:\\
Higher-rank composite-band alignment and hopping locality}}
\author{KSD reduced-Hamiltonian research program}
\date{Complete pre-execution D1 v1 protocol candidate}

\begin{document}
\maketitle

\begin{abstract}
This protocol defines a dimension-indexed synthetic family for testing how retained
rank, noncommuting gauge-circuit depth, admissible alignment, and spatial dimension
change matrix-valued hopping locality while the untruncated retained spectrum and
projector remain exactly equivalent. The common analytic construction permits arbitrary
declared finite rank and finite circuit depth and is topologically trivial by
construction. Its first bounded specialization, D1 P0/G0 v1, fixes a constant-projector
real-gauge model, exact exterior-isolation and alignment certificates, four ranks, four
depths, A0--A3 alignment channels, odd Fourier-refinement meshes, a disjoint evaluation
mesh, inversion-complete supports, numerical and convergence gates, claim-resolution
thresholds, strict schema identities, an independent direct-formula verifier, resource
ceilings, and retained-evidence policy. The protocol is complete but unimplemented and
unexecuted; it reports no new higher-rank numerical result and makes no material,
topological, or infinite-rank claim.
\end{abstract}

\noindent\textbf{Keywords:} composite bands; unitary gauge; operator alignment;
Fourier hopping; locality; truncation; convergence; retained evidence


\section{Status and authority}

This document defines a publication-independent experiment family for spatial dimensions
one, two, and three. The common mathematical family remains prospective outside the
selected one-dimensional specialization. The first P0/G0 specialization is now a
complete pre-execution protocol candidate with identity

\begin{quote}\small
\path{ksdft2effmass.campaign.periodic-reduction.periodic1d.higher-rank-composite-band-alignment-and-locality.v1}\par
\end{quote}

and short label \path{HR-CBAL-D1-P0-G0-v1}. Its model, cases, sample roles, tolerances,
supports, decision rules, schemas, verifier method, resource gate, and evidence policy
are selected below. The identity is reserved but is not promoted to an immutable
executed version until implementation, schema, and independent-verifier review pass.
No run or result is reported, and this document does not authorize execution.

The family is a successor to, not a replacement for, the retained rank-two controls and
the extracted fixed rank-three two-dimensional calculation. No rank-two tolerance, mesh,
gauge amplitude, truncation range, or evidence disposition transfers into this protocol.
The 2-D, 3-D, P1, and G1 branches remain prospective and cannot inherit the D1 v1
numerical controls without an explicit reviewed specialization.

\section{Scientific question}

For an isolated composite subspace of arbitrary finite rank, how do

\begin{enumerate}
\item retained rank,
\item noncommuting gauge-decomposition depth,
\item admissible alignment family, and
\item spatial dimension
\end{enumerate}

change the matrix-valued hopping locality and the spectral consequences of deliberate
finite-support truncation, even when the untruncated retained spectra and projectors are
exactly equivalent?

The protocol treats rank, decomposition depth, and spatial dimension as independent
variables. A calculation may vary only one while holding the other declared controls
fixed, or it must label the result as a joint variation. The phrase \textbf{arbitrary rank}
means that the formulas and software contract accept any declared finite integer rank
above the protocol minimum. It does not mean an infinite-dimensional computation,
uniform cost in rank, or numerical evidence for every rank.

\section{Common geometry and state spaces}

The first control uses the dimensionless orthogonal unit lattice in
$d\in\{1,2,3\}$ dimensions. Its reduced Brillouin torus is

\[
\mathbb T^d=[-\pi,\pi)^d,
\]

with periodic identification in every coordinate. Real-space translations are
$\mathbf R\in\mathbb Z^d$, measured with the Euclidean unit-lattice metric. A later
nonorthogonal or material specialization is a different protocol.

For retained rank $r\geq3$, the parent state space is $\mathbb C^{2r}$ with canonical
basis $\{|a\rangle\}_{a=1}^{2r}$. The embedding
$E_r:\mathbb C^r\rightarrow\mathbb C^{2r}$ maps retained coordinates to the first $r$
canonical coordinates, and

\[
P_0=E_rE_r^\dagger.
\]

Every result must record $d$, $r$, the ambient dimension $2r$, basis ordering, right-action
frame convention, dimensionless reciprocal coordinates, energy unit, scalar energy
reference, lattice metric, and Fourier-sign convention. Operators are never subtracted
across unequal ranks or unidentified spaces.

\section{Periodic scalar fields}

For a finite harmonic inventory $\mathcal Q_d\subset\mathbb Z^d\setminus\{0\}$, define

\[
\mathcal F_d(\mathbf k;c_0,a,b)
=c_0+\sum_{\mathbf q\in\mathcal Q_d}
\left[
 a_{\mathbf q}\sin(\mathbf q\cdot\mathbf k)
 +b_{\mathbf q}\cos(\mathbf q\cdot\mathbf k)
\right].
\]

The candidate inventories are

\[
\begin{aligned}
\mathcal Q_1={}&\{(1),(2)\},\\
\mathcal Q_2={}&\{(1,0),(0,1),(1,1),(1,-1)\},\\
\mathcal Q_3={}&\{(1,0,0),(0,1,0),(0,0,1),
(1,1,0),(0,1,1),(1,0,1),(1,1,1)\}.
\end{aligned}
\]

These sets specify candidate functional support, not numerical coefficients. A genuinely
multidimensional case must use at least one nonzero mixed-harmonic coefficient, and its
active wavevectors must span $\mathbb R^d$. Otherwise it is labeled separable or
effectively lower dimensional. The protocol assumes no point-group, time-reversal, or
inversion symmetry beyond constraints explicitly imposed by a reviewed coefficient
case.

\section{Isolated analytic parent}

Let

\[
D_r(\mathbf k)=\operatorname{diag}
\bigl(\epsilon_1(\mathbf k),\ldots,\epsilon_r(\mathbf k)\bigr)
\]

and

\[
D_{\mathrm{ext}}(\mathbf k)=\operatorname{diag}
\bigl(\epsilon^{\mathrm{ext}}_1(\mathbf k),\ldots,
\epsilon^{\mathrm{ext}}_r(\mathbf k)\bigr),
\]

where every dispersion is a real field from the declared $\mathcal F_d$ family. A case
is admissible only if it has the prospectively certified uniform exterior-isolation
energy

\[
\Delta_{\mathrm{iso}}=\inf_{\mathbf k\in\mathbb T^d}
\left[
\min_b\epsilon_b^{\mathrm{ext}}(\mathbf k)
-\max_a\epsilon_a(\mathbf k)
\right]>0.
\]

A finite sampled minimum is not this certificate. Until an analytic interval bound or a
reviewed verified enclosure is selected, a sampled isolation energy is reported only as
a sampled diagnostic. The symbol $E_g$ is reserved for a physical valence--conduction
band gap; $E_{\mathrm{ref}}$ is only the common positive energy scale.

For each retained coordinate define the anti-Hermitian retained--exterior generator

\[
K_a=|r+a\rangle\langle a|-|a\rangle\langle r+a|.
\]

With real periodic fields $\eta_a(\mathbf k)$, set

\[
W_r(\mathbf k)=\prod_{a=1}^{r}
\exp\!\left[\eta_a(\mathbf k)K_a\right],
\qquad
F_r(\mathbf k)=W_r(\mathbf k)E_r.
\]

The $K_a$ act on disjoint coordinate pairs, so the factors commute. The parent operator
and retained projector are

\[
H_p^{(d,r)}(\mathbf k)=W_r(\mathbf k)
\begin{pmatrix}
D_r(\mathbf k)&0\\
0&D_{\mathrm{ext}}(\mathbf k)
\end{pmatrix}
W_r^\dagger(\mathbf k),
\qquad
P_r(\mathbf k)=F_r(\mathbf k)F_r^\dagger(\mathbf k).
\]

The smooth represented operator is exactly

\[
H_{F_r}(\mathbf k)=F_r^\dagger H_p^{(d,r)}F_r=D_r(\mathbf k).
\]

Two projector controls remain separate:

\begin{itemize}
\item \textbf{P0, constant-projector control:} every $\eta_a=0$;
\item \textbf{P1, moving-projector control:} prospectively selected nonconstant $\eta_a$ fields.
\end{itemize}

P0 isolates internal gauge, alignment, Fourier conversion, and truncation. P1 additionally
tests frame construction, neighbor transport, reciprocal-boundary sewing, and projector
reconstruction. P0 and P1 results are not pooled into one error measure.

Internal retained-band crossings or degeneracies are permitted because the retained
projector and analytic frame are defined as complete rank-$r$ objects. Individual
backend-selected eigenvectors are not treated as canonical band labels.

\section{Layered Givens gauge decomposition}

For adjacent retained coordinates $j$ and $j+1$, define the embedded complex Givens
rotation

\[
R_{j,j+1}(\theta,\phi)=
\begin{pmatrix}
\cos\theta&-e^{-i\phi}\sin\theta\\
e^{i\phi}\sin\theta&\cos\theta
\end{pmatrix}_{j,j+1}.
\]

At layer $\ell$, use the disjoint neighbor pairs

\[
\mathcal P_\ell=
\left\{(j,j+1):1\leq j<r,\ j\equiv\ell\pmod 2\right\}.
\]

Let every $\theta_{\ell j}(\mathbf k)$, $\phi_{\ell j}(\mathbf k)$, and
$\chi_a(\mathbf k)$ be a declared periodic field. For finite circuit depth $L\geq1$,
define

\[
Q_r(\mathbf k)=\operatorname{diag}
\left(e^{i\chi_1(\mathbf k)},\ldots,e^{i\chi_r(\mathbf k)}\right)
\]

and

\[
G_{r,L}(\mathbf k)=Q_r(\mathbf k)
\prod_{\ell=1}^{L}
\prod_{(j,j+1)\in\mathcal P_\ell}
R_{j,j+1}\!\left(
\theta_{\ell j}(\mathbf k),\phi_{\ell j}(\mathbf k)
\right).
\]

Increasing layer order is part of the canonical definition. Factors within a layer
commute; adjacent layers generally do not. The interaction graph is connected once the
selected layers cover every adjacent pair, so the construction is not a direct sum of
independent rank-two attacks.

The protocol distinguishes two gauge classes:

\begin{itemize}
\item \textbf{G0, real class:} $\phi_{\ell j}=0$ and $\chi_a=0$, hence
$G_{r,L}\in SO(r)$;
\item \textbf{G1, complex class:} complex pair phases and diagonal phases are allowed, hence
$G_{r,L}\in U(r)$.
\end{itemize}

A G0 result does not validate G1. The attacked frame and represented operator are

\[
F'_{r,L}(\mathbf k)=F_r(\mathbf k)G_{r,L}(\mathbf k),
\qquad
H'_{r,L}(\mathbf k)=G_{r,L}^\dagger(\mathbf k)
D_r(\mathbf k)G_{r,L}(\mathbf k).
\]

The defining exact identities are

\[
F_{r,L}^{\prime\dagger}F'_{r,L}=I_r,
\qquad
F'_{r,L}F_{r,L}^{\prime\dagger}=P_r,
\qquad
\operatorname{spec}H'_{r,L}=\operatorname{spec}D_r.
\]

These are model identities and must be tested to numerical roundoff after sampling. They
are not empirical scientific conclusions. Increasing $L$ changes a representation of the
same retained operator unless another case explicitly changes the parent.

\section{Alignment families}

Every comparison declares its admissible alignment family before observing results.
The initial diagnostic ladder is:

\begin{enumerate}
\item \textbf{A0, identity:} no alignment;
\item \textbf{A1, one global unitary:} one $U\in U(r)$ for the complete sampled torus;
\item \textbf{A2, pointwise complete unitary:} one independently selected full-rank
$U(\mathbf k)\in U(r)$ at each construction point;
\item \textbf{A3, known inverse oracle:} the analytic $G_{r,L}^{-1}(\mathbf k)$, used only to
verify construction and convention;
\item \textbf{A4, restricted circuit:} a separately frozen depth- and harmonic-bounded Givens
family, if a reviewed optimizer and certificate policy are later selected.
\end{enumerate}

A1, A2, and A4 are different feasible sets. None is described as a numerical
approximation to another. The oracle A3 is not an admissible fitted solution and cannot
be used to claim practical recoverability.

Pointwise and aggregate Procrustes selection must disclose absolute and relative
singular-value cutoffs. Equality at the applied cutoff is deficient. A rank-deficient
cross-covariance fails closed; no null-space completion is inferred. The retained output
must include the complete singular-value inventory or sufficient lossless summaries to
verify every rank disposition independently.

Matrices are aligned before Fourier transformation. Aligning already truncated hopping
blocks is a different, generally noncommuting operation and is outside this protocol.

\section{Boundary to algorithm-selected localization}

D1 v1 supplies every frame analytically and does not execute Wannier90 or optimize a
spread functional. This exclusion is deliberate: it keeps parent construction, internal
gauge attack, admissible alignment, Fourier conversion, and truncation independently
identifiable. The retained one-dimensional and two-dimensional Wannier90 calculations
reported by Paper 3 are historical algorithm-selected controls, not runs of this
protocol.

Initialization, native convergence, preconditioning, inactive embedding, optimizer
continuation, and observed-basin reproducibility belong to a separate
Wannier90-localization-reproducibility experiment family. Evidence from that family may
qualify claims about algorithm-selected frames but cannot select or modify a D1 v1 mesh,
tolerance, support, or disposition.

\section{Selected v1 candidate: periodic 1-D real gauge}

The first bounded specialization is \path{HR-CBAL-D1-P0-G0-v1}, with the full identity given
above. It isolates retained rank and Givens-circuit depth before projector motion,
complex phases, or multidimensional plaquettes are introduced. All numerical execution
uses binary64 arithmetic and sets $E_{\mathrm{ref}}=1$ in the declared dimensionless
energy unit; reported dimensional formulas retain $E_{\mathrm{ref}}$ explicitly.

\subsection{Rank-normalized parent}

For $a\in\{1,\ldots,r\}$ define the rank-normalized orbital coordinate

\[
x_{r,a}=\frac{2a-r-1}{r-1}\in[-1,1]
\]

and the bounded even harmonic

\[
c(k)=\frac34\cos k+\frac14\cos(2k),
\qquad |c(k)|\leq1.
\]

In units of the declared positive reference energy $E_{\mathrm{ref}}$, the retained and
exterior dispersions are

\[
\frac{\epsilon_a(k)}{E_{\mathrm{ref}}}
=0.40x_{r,a}+0.25\left(1+0.20x_{r,a}\right)c(k),
\]

\[
\frac{\epsilon_a^{\mathrm{ext}}(k)}{E_{\mathrm{ref}}}
=3.00+0.30x_{r,a}+0.10c(k).
\]

For fixed $k$, the retained maximum occurs at $x_{r,a}=1$ and the exterior minimum at
$x_{r,a}=-1$. Their separation is

\[
\frac{\min_a\epsilon_a^{\mathrm{ext}}(k)-\max_a\epsilon_a(k)}
{E_{\mathrm{ref}}}
=2.30-0.20c(k)\geq2.10.
\]

Because $c(0)=1$, equality is attained. Every selected rank therefore has the exact,
mesh-independent isolation certificate
$\Delta_{\mathrm{iso}}=2.10E_{\mathrm{ref}}$.

The first tranche uses P0: $\eta_a(k)=0$, $W_r(k)=I_{2r}$, and
$F_r(k)=E_r$. The parent and retained projector are consequently known without an
eigensolver. Internal retained-band crossings, if present, do not change the complete
rank-$r$ projector.

\subsection{Deterministic real attack}

For $j\in\{1,\ldots,r-1\}$ define

\[
y_{r,j}=\frac{2j-r}{r-1},
\qquad |y_{r,j}|<1,
\]

and, for each positive layer index $\ell$,

\[
s_\ell(k)=\frac34\sin k+\frac14(-1)^\ell\sin(2k).
\]

The selected adjacent-pair angle is

\[
\theta_{\ell j}(k)=0.35(-1)^{\ell+j}
\left(1+0.15y_{r,j}\right)s_\ell(k).
\]

All pair phases and diagonal phases vanish, so
$G_{r,L}(k)\in SO(r)$. To state the exact continuum angle maximum, define

\[
q=\frac{\sqrt{41}-3}{8},
\qquad
s_*=\sqrt{1-q^2}\left(\frac34+\frac q2\right).
\]

For either sign in $s_\ell$, maximizing over $k$ gives $\max_k|s_\ell(k)|=s_*$;
the maximizing cosine is $q$ for even $\ell$ and $-q$ for odd $\ell$. The positive
amplitude factor is largest at $j=r-1$. Hence

\[
\max_{\substack{k\in[-\pi,\pi)\\ \ell\geq1,\ 1\leq j<r}}
|\theta_{\ell j}(k)|
=0.35\left[1+0.15\frac{r-2}{r-1}\right]s_*<0.4025\ \mathrm{rad}.
\]

This is an analytic continuum certificate, unlike the sampled development diagnostic
below. Layers use the canonical odd/even pair sets and increasing product order defined
above. One layer is an intentionally incomplete-connectivity baseline; two or more
layers cover every adjacent edge and generate the noncommuting connected-chain control.

\subsection{Exact representation certificates}

The smooth represented field $D_r(k)$ is a trigonometric polynomial with harmonics only
at $0$, $1$, and $2$. Under the continuum Fourier convention of this protocol, its
blocks therefore satisfy

\[
T_{D_r}(R)=0\qquad\text{for every }|R|>2.
\]

This is an analytic support statement, not a sampled locality result. By contrast,
entries of $G_{r,L}(k)$ contain sines and cosines of nonconstant trigonometric
polynomials. Consequently the attacked field
$H'_{r,L}(k)=G_{r,L}^{\dagger}(k)D_r(k)G_{r,L}(k)$ is periodic and analytic but is not
assigned any finite Fourier support by the construction. The first-two-harmonic input
inventory must not be mistaken for a hopping cutoff after the nonlinear gauge attack.

For P0, the exact candidate-to-reference cross-overlap is

\[
M(k)=F_{r,L}^{\prime\dagger}(k)F_r(k)=G_{r,L}^{\dagger}(k).
\]

Every singular value of $M(k)$ is exactly one. The unique complete-unitary pointwise
polar factor is therefore $G_{r,L}^{\dagger}(k)=G_{r,L}^{-1}(k)$. In exact arithmetic,
A2 and A3 both restore $D_r(k)$ before Fourier transformation and hence restore its
$|R|\leq2$ support. They remain separate diagnostic channels: A2 must reconstruct the
polar factor from the sampled frames, whereas A3 receives the generating inverse. A
finite-grid agreement between them tests implementation and convention only; it is not
new scientific evidence. No corresponding recovery statement applies to A1, whose
aggregate cross-covariance must independently pass its rank threshold.

Every represented field is pointwise Hermitian and periodic, so its continuum Fourier
blocks obey $T(-R)=T(R)^\dagger$. A finite transform must test the corresponding discrete
pairing rather than assume it.

\subsection{Rank and depth ladder}

The selected full one-dimensional development grid is

\[
r\in\{3,4,8,16\},
\qquad
L\in\{1,2,4,8\}.
\]

The shared cross-dimensional core for later 2-D and 3-D specializations is the subset
$r\in\{3,4,8\}$ and $L\in\{1,2,4\}$. Selecting that common subset now does not select
2-D or 3-D harmonic coefficients, meshes, or evidence status.

Every case evaluates A0 identity, A1 one-global-unitary, A2 pointwise-complete-unitary,
and A3 known-inverse-oracle channels. Canonical case identifiers are
\texttt{r\{rank:02d\}-L\{depth:02d\}} and channel identifiers are exactly \path{A0}, \path{A1}, \path{A2}, and
\path{A3}; the schemas enumerate the resulting 16 case identifiers rather than accept an
arbitrary string. A4 restricted-circuit fitting is excluded from this tranche because no
optimizer or certificate policy has been selected. A3 remains a construction oracle and
cannot be reported as a fitted recovery result.

For cross-rank reporting, matrix Frobenius quantities must be retained both raw and
normalized by $\sqrt rE_{\mathrm{ref}}$; quadratic matrix weights must be retained both
raw and normalized by $rE_{\mathrm{ref}}^2$. This normalization is selected before
observing outcomes and does not replace the raw values. Ordered-spectrum maximum errors
are normalized by $E_{\mathrm{ref}}$, and spectrum RMS errors by
$\sqrt rE_{\mathrm{ref}}$.

\subsection{Frozen numerical controls}

The three odd, centered, half-open construction meshes are

\[
\mathcal K_M=\left\{-\pi+\frac{2\pi n}{M}:n=0,\ldots,M-1\right\},
\qquad M\in\{513,1025,2049\}.
\]

The $M=2049$ mesh is the final construction mesh. The smaller meshes are mandatory
refinement controls, not optional pilot runs. A1 is selected once from the normalized
aggregate cross-covariance on the final construction mesh,

\[
B_{2049}=\frac1{2049}\sum_{k\in\mathcal K_{2049}}
F_{r,L}^{\prime\dagger}(k)F_r(k),
\]

and that one constant unitary is applied to every Fourier-refinement mesh and to direct
evaluation. A2 is reconstructed independently at each construction point. A3 evaluates
the analytic inverse and is never fitted.

The withheld evaluation mesh is

\[
\mathcal K_{\mathrm{eval}}=
\left\{-\pi+\frac{2\pi(m+1/2)}{8191}:m=0,\ldots,8190\right\}.
\]

It is disjoint from every construction mesh: equality would require an even integer to
equal an odd integer after clearing the odd denominators. Evaluation cannot change any
alignment, support, tolerance, case, or reporting rule.

For each odd $M$, the complete discrete translation inventory is
$R=-(M-1)/2,\ldots,(M-1)/2$. The prospectively fixed truncation supports are

\[
\Omega_S=\{R\in\mathbb Z:|R|\leq S\},
\qquad S\in\{0,1,2,4,8,16,32,64,128\}.
\]

Every support contains complete inversion pairs. No first-passing support is selected;
all nine supports are reported for every applicable channel, rank, depth, and mesh.

A singular-value decision uses

\[
\tau_{\mathrm{sv}}(X)=
\max\!\left(10^{-12},10^{-10}\sigma_{\max}(X)\right).
\]

A cross-covariance is full rank only when
$\sigma_{\min}(X)>\tau_{\mathrm{sv}}(X)$; equality is deficient. A deficient A1
aggregate makes A1 downstream diagnostics inapplicable for that case but does not erase
A0, A2, or A3. A deficient A2 point makes the complete case invalid because P0 supplies
an analytically full-rank overlap.

All pass comparisons are inclusive except the strict singular-value test above. Norms
are Frobenius unless another norm is named. A normalized frame or unitary defect is the
raw Frobenius norm divided by $\sqrt r$; a normalized operator or hopping-block defect is
divided by $\sqrt rE_{\mathrm{ref}}$. The relative Parseval defect divides the absolute
difference by the larger of the two compared quadratic weights and
$rE_{\mathrm{ref}}^2\times10^{-30}$. The frozen numerical tolerances are:

\begin{center}\small
\begin{tabularx}{\textwidth}{YYY}
\toprule
\textbf{Diagnostic} & \textbf{Normalization} & \textbf{Tolerance} \\
\midrule
momentum and endpoint coordinate construction & radians & $10^{-14}$ \\
frame orthonormality, gauge unitarity, projector idempotency & raw Frobenius norm & $10^{-12}$ \\
smooth/attacked projector equality & raw Frobenius norm & $10^{-12}$ \\
direct ordered-spectrum equivalence & $E_{\mathrm{ref}}$ & $10^{-11}$ \\
A2 and A3 frame recovery & $\sqrt r$ & $10^{-11}$ \\
A2 and A3 represented-operator recovery & $\sqrt rE_{\mathrm{ref}}$ & $10^{-11}$ \\
reciprocal-endpoint sewing and primitive closure & normalized Frobenius norm & $10^{-11}$ \\
maximum closure-eigenphase magnitude & radians & $10^{-10}$ \\
Fourier reconstruction on a construction mesh & $\sqrt rE_{\mathrm{ref}}$ & $2\times10^{-11}$ \\
Fourier Hermitian-pairing defect & $\sqrt rE_{\mathrm{ref}}$ & $2\times10^{-11}$ \\
relative Parseval defect & dimensionless & $2\times10^{-11}$ \\
final full-interpolant operator error on evaluation mesh & $\sqrt rE_{\mathrm{ref}}$ & $10^{-9}$ \\
final full-interpolant ordered-spectrum maximum error & $E_{\mathrm{ref}}$ & $10^{-10}$ \\
retained-block change from $M=1025$ to $M=2049$, $\lvert R\rvert\leq128$ & $\sqrt rE_{\mathrm{ref}}$ per block & $10^{-10}$ \\
\bottomrule
\end{tabularx}
\end{center}

The odd meshes avoid a self-inverse Nyquist block, each refinement approximately doubles
resolution, and even the coarsest mesh leaves translations beyond the largest retained
support. The evaluation mesh is roughly four times the final construction size. The
$10^{-9}$--$10^{-10}$ evaluation gates sit well above binary64 roundoff while remaining
at least four orders of magnitude below the $10^{-6}$ effect-resolution floor. These are
prospective resolution choices, not tolerances fitted to an observed outcome.

For a nonnegative scalar locality diagnostic $z$, refinement stability between
$M=1025$ and $M=2049$ requires

\[
|z_{2049}-z_{1025}|\leq10^{-10}+10^{-6}
\max(|z_{2049}|,|z_{1025}|)
\]

after the diagnostic's declared rank and energy normalization. All cases and applicable
channels must satisfy the construction, equivalence, conversion, evaluation, and
refinement gates. A failed gate invalidates scientific interpretation of that case; it
is not repaired by relaxing a tolerance after inspection.

\subsection{Non-authoritative finite-grid sanity check}

The development utility \path{verify_d1_p0_g0_controls.py} samples the selected formulas on

\[
\mathcal K_M=\left\{-\pi+\frac{2\pi n}{M}:n=0,\ldots,M\right\},
\qquad M=20000,
\]

and reports, for each selected rank,

\[
\widehat\Delta_{\mathrm{iso},r}
=\min_{k\in\mathcal K_M}
\left[\min_a\epsilon_a^{\mathrm{ext}}(k)-\max_a\epsilon_a(k)\right]
\]

and

\[
\widehat\theta_r
=\max_{\substack{k\in\mathcal K_M,\ 1\leq\ell\leq8\\1\leq j<r}}
|\theta_{\ell j}(k)|.
\]

This grid is not a selected construction mesh or evaluation mesh. The sampled isolation
energy and sampled angle maximum are only reproducibility checks against the analytic
certificates. The sampled angle is slightly below the exact continuum maximum because
the maximizing momentum need not lie on the grid. Neither sampled value may be promoted
to convergence or retained-evidence status.

\section{Dimension-specific transport and sewing}

All analytic fields are $2\pi$ periodic, so the canonical ambient basis has identity
face sewing. A numerical route must nevertheless reconstruct and verify the declared
sewing rather than infer it from array shape.

\subsection{Periodic 1-D}

The one-dimensional case has one neighbor-link family, one primitive reciprocal closure,
and hopping blocks indexed by $R\in\mathbb Z$. It is the primary rank/depth isolation
case because it introduces no plaquette channel. Closure eigenphases are stored as an
unordered rank-$r$ multiset.

\subsection{Periodic 2-D}

The two-dimensional case has $x$- and $y$-link families, sewing across both reciprocal
face pairs, primitive closures in both directions, and $xy$ plaquettes. Hopping blocks
are indexed by $\mathbf R\in\mathbb Z^2$. Wilson-loop phase spectra are compared as
unordered multisets; no band-by-band phase labeling is inferred through crossings.

\subsection{Periodic 3-D}

The three-dimensional case has $x$, $y$, and $z$ link families, sewing across all three
reciprocal-face pairs, three primitive closures, and $xy$, $xz$, and $yz$ plaquettes.
Hopping blocks are indexed by $\mathbf R\in\mathbb Z^3$. Dense tensor fields must not be
materialized solely for diagnostics when a matrix-free or streamed route can provide the
same declared quantity.

Neighbor transport is applicable only if every required cross-overlap passes its applied
singular-value threshold. Closure and plaquette diagnostics are inapplicable, rather than
zero, when a required link is rank deficient.

\section{Topological boundary}

The analytic frame $F_r(\mathbf k)=W_r(\mathbf k)E_r$ is global and periodic. Therefore
P0 and P1 define topologically trivial complex composite bundles. Wilson loops, closure
phases, and plaquettes are transport and sewing diagnostics in this family; they are not
claims of Chern, $\mathbb Z_2$, crystalline, or three-dimensional topological character.

A nonidentity Wilson loop does not by itself imply a topological obstruction. A future
obstructed model would require patchwise frames and transition functions, and would be a
separate experiment family rather than a hidden coefficient choice here.

The reusable diagnostics developed from this family may later support a fail-closed
pristine/defect compatibility screen, but Paper 3 does not construct defect operators.
Paper 5 owns the common-space subtraction decision, and Paper 6 owns full topological
and symmetry-obstruction protocols. The synthetic P0/P1 cases exercise known-compatible
frames; they do not validate material bundle compatibility.

\section{Fourier representation and locality}

For a declared represented field $H(\mathbf k)$, define the continuum hopping blocks

\[
T(\mathbf R)=\frac{1}{(2\pi)^d}
\int_{\mathbb T^d}e^{-i\mathbf k\cdot\mathbf R}
H(\mathbf k)\,d^d\mathbf k.
\]

For D1 v1, the selected discrete convention on $\mathcal K_M$ is

\[
T_M(R)=\frac1M\sum_{k\in\mathcal K_M}e^{-ikR}H(k),
\qquad
H_M(k)=\sum_{R=-(M-1)/2}^{(M-1)/2}e^{ikR}T_M(R).
\]

Only the forward transform carries $1/M$. The odd mesh gives a unique inversion-symmetric
canonical translation inventory with no Nyquist singleton. The discrete transform is a
finite representation of sampled data, not the continuum coefficient sequence.
Continuum quadrature error, finite-mesh aliasing, and deliberate hopping truncation
remain distinct.

Every support is an explicit set of translation vectors. It must include inversion pairs
and complete orbits of any point group actually declared by the case. Equal-radius ties
are never split. A shell number is not transferred between dimensions or metrics without
the full vector inventory and physical radius.

For a retained support $\Omega\subset\mathbb Z^d$, define

\[
T_\Omega(\mathbf R)=
\begin{cases}
T(\mathbf R),&\mathbf R\in\Omega,\\
0,&\mathbf R\notin\Omega.
\end{cases}
\]

The protocol reports complete matrix-block norms, omitted quadratic weight, omitted block
norm, reconstruction residual, and the spectral consequence of the reconstructed
truncated operator. Frobenius quantities are reported both unnormalized and, when
cross-rank comparison is intended, with a prospectively selected rank normalization.
The normalization and its interpretation must be frozen before results; dividing by
$\sqrt r$ after observing a trend is not permitted.

\section{Construction and evaluation procedure}

For each of the 16 $(r,L)$ cases, the producer performs the following fixed sequence.

\begin{enumerate}
\item Reconstruct $D_r(k)$, the exterior block, $G_{r,L}(k)$, $F_r(k)$, and
$F'_{r,L}(k)$ directly from the canonical formulas on all three construction meshes.
\item Evaluate analytic endpoint values at both $-\pi$ and $+\pi$ for sewing; a duplicated
endpoint is never inserted into a Fourier mesh.
\item Construct neighbor links in increasing mesh-index order with the last point linked to
the first through the declared identity face sewing. Store the minimum threshold
margin, maximum unitary defect, closure defect, and unordered closure eigenphases.
\item Select A1 once on $\mathcal K_{2049}$. Construct A2 independently on every mesh and
evaluate A3 from the known inverse. Record A1 as deficient rather than inventing a
null-space completion when its aggregate fails the strict threshold.
\item Align matrices before Fourier transformation and compute complete discrete blocks for
every applicable channel on every construction mesh.
\item Evaluate all fixed supports without selecting a preferred support. Reconstruct the
full and truncated interpolants on the withheld mesh.
\item Compare the full interpolant with the applicable direct analytic evaluation field:
$H'_{r,L}$ for A0, $U_{\mathrm{glob}}^\dagger H'_{r,L}U_{\mathrm{glob}}$ for a
qualified A1, and $D_r$ for A2 and A3.
\item Apply every integrity, refinement, and claim-resolution rule without modifying any
control.
\end{enumerate}

The evaluation role is withheld rather than called training: no statistical model is
learned. Evaluation data cannot alter ranks, depths, coefficients, meshes, alignments,
supports, thresholds, output selection, or figure selection.

\section{Diagnostic definitions and ledger}

For a channel $A$ and construction mesh $M$, write $T_{A,M}(R)$ for the complete block
and $H_{A,M}(k)$ for its full interpolant. For support $\Omega_S$, retain both

\[
W_{A,M}(S)=\sum_{R\notin\Omega_S}\|T_{A,M}(R)\|_F^2,
\qquad
N_{A,M}(S)=\sqrt{W_{A,M}(S)},
\]

and the block-$\ell^1$ tail

\[
L_{A,M}(S)=\sum_{R\notin\Omega_S}\|T_{A,M}(R)\|_F.
\]

Report $W$ raw and divided by $rE_{\mathrm{ref}}^2$; report $N$ and $L$ raw and divided
by $\sqrt rE_{\mathrm{ref}}$. These are complete-matrix quantities, not sums of selected
orbital entries.

On the withheld mesh, keep the following operator and spectral channels separate:

\[
\begin{aligned}
e_{\mathrm{interp}}(k)&=H_{A,2049}(k)-H_A^{\mathrm{direct}}(k),\\
e_{\mathrm{trunc}}^{(S)}(k)&=H_{A,2049}^{(S)}(k)-H_{A,2049}(k),\\
e_{\mathrm{total}}^{(S)}(k)&=H_{A,2049}^{(S)}(k)-H_A^{\mathrm{direct}}(k).
\end{aligned}
\]

For each, report maximum and RMS Frobenius operator errors and maximum and RMS ordered-
eigenvalue errors with the fixed cross-rank normalizations. The three channels are not
added, and no cancellation between interpolation and truncation is called convergence.

Every case retains separate diagnostics for:

\begin{enumerate}
\item parent Hermiticity, basis order, scalar energy reference, and exact analytic gap;
\item sampled retained/exterior separation on every mesh, labeled sampled;
\item frame orthonormality, gauge unitarity, projector idempotency, and projector equality;
\item direct ordered-spectrum equality before and after each applicable alignment;
\item A1 aggregate singular values and threshold margin;
\item for A2, minimum singular value, maximum singular value, minimum threshold margin,
deficient-point count, and the exact indices of any deficient points;
\item neighbor-link singular-value extrema and deficient-link indices;
\item reciprocal-endpoint sewing, primitive closure, and unordered closure phases;
\item A0--A3 frame and represented-operator residuals, with inapplicable distinguished from
numerical zero;
\item Fourier reconstruction, Hermitian pairing, Parseval, and canonical-index checks;
\item every block norm for $|R|\leq128$ and the three tail quantities for every support;
\item withheld interpolation error;
\item deliberate truncation error relative to the full interpolant; and
\item total error relative to the direct analytic evaluation field.
\end{enumerate}

Parent construction, covariant equivalence, restricted alignment, Fourier conversion,
interpolation, and deliberate truncation remain distinct. No combined diagnostic score
is defined.

\section{Convergence and claim-resolution policy}

A D1 v1 case is \textbf{numerically qualified} only if every mandatory integrity gate passes,
its final interpolation errors pass, and all retained blocks and locality diagnostics
meet the frozen $1025\rightarrow2049$ refinement rules. A1 may instead receive the
explicit \path{aggregate-rank-deficient} disposition; this is a bounded A1 result, not a
failed search and not permission to complete the map.

The exact support and A2/A3 recovery statements are analytic certificates and never
become empirical findings. Numerical values only test their implementation.

For a qualified case, define a \textbf{resolved gauge-induced locality extension at $S=2$}
only when

\[
\frac{N_{A0,2049}(2)}{\sqrt rE_{\mathrm{ref}}}\geq10^{-6},
\qquad
\frac{N_{A3,2049}(2)}{\sqrt rE_{\mathrm{ref}}}\leq10^{-10},
\]

and the A0 total ordered-spectrum maximum error after $S=2$ truncation is at least
$10^{-6}E_{\mathrm{ref}}$. Below-threshold behavior is reported as unresolved at v1
resolution, not as proof of equal locality. A qualified A1 is classified as a
\textbf{resolved global-alignment residual} only if its maximum normalized represented-
operator residual is at least $10^{-6}$; otherwise it is unresolved at v1 resolution.
No ratio with a roundoff-scale denominator is used.

Rank or depth differences are descriptive unless the absolute difference between two
normalized diagnostics exceeds both $10^{-6}$ and ten times their larger measured
refinement change. Monotonicity is not assumed, and the finite ladder supports no
asymptotic-in-rank or asymptotic-in-depth law. A failure to resolve an effect is not a
negative theorem.

\section{Scaling study}

The selected v1 matrix contains four rank sweeps at every fixed depth and four depth
sweeps at every fixed rank. Rank comparisons hold dimension, parent coefficient rule,
active harmonics, gauge class, supports, and $E_{\mathrm{ref}}$ fixed. Depth comparisons
hold rank and the per-layer rule fixed. The shared later-dimensional subset creates no
D1-to-D2 or D1-to-D3 claim until those protocols select matched meshes and metrics.

Dense unitary factorization scales as $O(r^3)$ per sampled point. Actual peak memory,
wall time, and processed mesh points are recorded per producer and verifier invocation.
Resource failure is unresolved coverage, not a scientific failure of the model.

\section{Independent verification}

The selected verifier method identity is

\begin{quote}\small
\path{ksdft2effmass.verify.periodic-reduction.periodic1d.higher-rank-composite-band-alignment-and-locality.direct-formula.v1}\par
\end{quote}

The verifier reconstructs all inputs from \path{input.json} and the canonical protocol; it
does not trust producer-derived arrays. It independently composes scalar fields, Givens
blocks, frames, parents, projectors, alignment objectives, neighbor links, closures,
Fourier sums, supports, interpolants, and every error channel. It uses explicit direct
Fourier sums for retained blocks and does not call the producer experiment action, its
FFT wrapper, its serializer, or a reusable campaign workflow action.

Shared Python, NumPy, BLAS, and LAPACK lineage is permitted only when disclosed in
provenance. The verifier uses its own polar/SVD call site and independently checks the
objective value, singular values, threshold margin, and unitary defect. It never silently
completes a deficient map. Producer and verifier results must agree within
$10^{-11}$ for normalized scalar diagnostics and within
$10^{-10}+10^{-9}\max(|x|,|y|)$ for normalized locality diagnostics. Every categorical
disposition must agree exactly.

Passing software tests or independent reconstruction establishes bounded software and
finite-protocol numerical behavior only. It does not establish material relevance,
localization optimality, nontrivial topology, continuum behavior outside the selected
analytic certificates, or validity at ranks and depths not executed.

\section{Schemas and canonical serialization}

The implementation must provide strict JSON Schemas with these immutable identifiers:

\begin{quote}\small
\path{urn:ksdft2effmass:campaign:periodic-reduction:periodic1d:higher-rank-composite-band-alignment-and-locality:v1:input}\par
\path{urn:ksdft2effmass:campaign:periodic-reduction:periodic1d:higher-rank-composite-band-alignment-and-locality:v1:result}\par
\path{urn:ksdft2effmass:campaign:periodic-reduction:periodic1d:higher-rank-composite-band-alignment-and-locality:v1:verification}\par
\path{urn:ksdft2effmass:campaign:periodic-reduction:periodic1d:higher-rank-composite-band-alignment-and-locality:v1:provenance}\par
\end{quote}

Every object rejects unknown properties. Numeric fields reject booleans, strings,
nonfinite values, and implicit units. Enumerations bind the exact ranks, depths, meshes,
supports, channels, conventions, status values, and method identities. \path{input.json}
contains the complete model coefficients and controls rather than only a protocol path.
\path{result.json} contains every case disposition, raw and normalized diagnostic, support
inventory, refinement comparison, resource measurement, and claim-resolution outcome.
\path{verification.json} contains independently reconstructed values and per-field comparison
dispositions. \path{provenance.json} binds source and runtime identities.

Canonical JSON is UTF-8, key-sorted, compact, finite, and terminated by exactly one line
feed. The implementation uses separators \path{(',', ':')}, preserves JSON array order where
order is scientific, and forbids NaN and infinity. SHA-256 is computed over canonical
bytes. Schema files themselves are checksummed and bound into provenance.

\section{Resource and execution gate}

D1 v1 is a deterministic local synthetic calculation. It must not invoke Quantum
ESPRESSO, Wannier90, a scheduler, a remote service, a GPU, or network access. The
producer and verifier each run as one process with numerical-library thread counts fixed
to one. The prospective ceilings are:

\begin{center}\small
\begin{tabularx}{\textwidth}{YYY}
\toprule
\textbf{Resource} & \textbf{Producer ceiling} & \textbf{Verifier ceiling} \\
\midrule
wall time & 30 minutes & 45 minutes \\
peak resident memory & 4 GiB & 4 GiB \\
retained package size, excluding generated figure images & 50 MiB total & included \\
\bottomrule
\end{tabularx}
\end{center}

The implementation must stream across rank/depth/channel cases and must not retain a
complete all-case mesh-by-matrix tensor solely for convenience. Exceeding a ceiling gives
\path{resource-limit-exceeded}, retains no successful evidence package, and triggers protocol
review rather than an unreviewed increase. Runtime, peak resident memory, thread
environment, platform, Python, NumPy, BLAS/LAPACK, and floating-point identifiers are
recorded.

Implementation review must pass unit tests, strict schema validation, type checking,
linting, source-manifest reconstruction, and an unreleased end-to-end dry run before any
retained execution is authorized. A dirty worktree is not represented as a clean Git
revision; all in-scope dirty source bytes must be included in \path{source.sha256}.

\section{Retained-evidence and failure policy}

A run identity has the no-overwrite form
\path{<version-id>.run-YYYYMMDDTHHMMSSZ} in UTC; a timestamp collision is an error rather than
permission to replace an existing directory. A successful retained package contains:

\begin{itemize}
\item canonical \path{input.json}, \path{result.json}, \path{verification.json}, and \path{provenance.json};
\item the four exact schemas used for validation;
\item \path{diagnostics.csv} with one explicitly keyed row per case, mesh, channel, and support;
\item \path{report.md} and publication-view figure data;
\item \path{source.sha256}, \path{SHA256SUMS}, and \path{evidence-package-identity.txt}; and
\item generated figures only as non-authoritative views of checksummed numerical records.
\end{itemize}

\path{SHA256SUMS} covers every retained file except itself and
\path{evidence-package-identity.txt}. The evidence identity is
\path{sha256:<SHA-256(SHA256SUMS bytes)>}. Consumers verify digests and schemas before reading
scientific values. The package binds the exact Git base revision, dirty-source manifest,
dependency revisions and lock digest, protocol and schema digests, run identity, verifier
method, runtime environment, and measured resources.

A successful package requires all 16 cases, all three construction meshes, all nine
supports, A0/A2/A3, every qualified A1 or exact A1 deficiency disposition, and complete
independent-verifier agreement. Partial rank/depth coverage, missing refinement levels,
failed schemas, failed digests, failed numerical gates, or resource-limit failures cannot
be promoted as successful evidence. Failed attempts may be kept only under a distinct
failed-run identity with incomplete outputs and the failure made explicit. They are not
publication evidence and never replace a successful package.

Retained evidence supports only the selected finite synthetic D1 cases and the declared
resolution classifications. It does not validate a material, a global localization
optimum, an unrestricted alignment algorithm, arbitrary executed rank, a continuum decay
law, or any topological classification. Publication, release, and evidence promotion
remain separate human decisions.

\section{Selections remaining outside D1 v1}

No D1 v1 scientific or numerical control remains outcome-selectable. Implementation
class names, file layout, and exact dependency revisions may be chosen during
implementation review only when they preserve every contract above; the final revisions
and bytes are then bound in provenance.

The following remain open for later, separately versioned protocols:

\begin{enumerate}
\item P1 moving-projector coefficients and their certified exterior-gap method;
\item every 2-D and 3-D coefficient table, mesh, support, and dimensional identity;
\item G1 complex phases and diagonal-phase laws;
\item A4 restricted-circuit optimizer, initialization, stopping rule, and certificate;
\item nonorthogonal metrics, point-group constraints, or material specializations; and
\item any topologically obstructed or patchwise-frame family.
\end{enumerate}

None of these may be added to D1 v1 after observing its outcomes. The D1 identity,
coefficients, exact certificates, rank/depth ladder, meshes, supports, tolerances,
A0--A3 channels, decision rules, verifier method, resource gate, and evidence policy are
now selected. Until reviewed implementation and schemas exist, the version remains
unexecuted and cannot supply higher-rank numerical evidence.
\end{document}
